\section{Công nghệ TinyML trên vi điều khiển}

TinyML (Tiny Machine Learning) là hướng tiếp cận đưa các mô hình học máy xuống thực thi trực tiếp trên các hệ thống nhúng và vi điều khiển với mức tiêu thụ điện năng chỉ vài miliwatt \cite{ref_1_1}. Việc tích hợp xử lý tín hiệu số (DSP) và mạng nơ-ron sâu ngay trên các thiết bị có tài nguyên hạn chế giúp dịch chuyển năng lực tính toán ra sát nguồn phát sinh dữ liệu. Phương pháp này giúp khắc phục hiệu quả các hạn chế về độ trễ, băng thông truyền tải và rủi ro bảo mật dữ liệu so với mô hình điện toán đám mây truyền thống \cite{ref_1_4}.

\subsection{Ràng buộc về tài nguyên phần cứng}

Triển khai mạng nơ-ron trên thiết bị biên như Seeed Studio XIAO ESP32-S3 yêu cầu tính toán kỹ lưỡng các giới hạn vật lý của phần cứng. Khác với môi trường đám mây, vi điều khiển phải đối mặt với ba rào cản thiết kế chính. Thứ nhất là giới hạn nghiêm ngặt về bộ nhớ, bao gồm bộ nhớ Flash để lưu trữ trọng số và bộ nhớ SRAM dành cho việc cấp phát tensor. Mặc dù ESP32-S3 sở hữu vi xử lý lõi kép Xtensa LX7 mạnh mẽ, dung lượng SRAM nội (internal SRAM) chỉ ở mức 512 KB \cite{ref_1_15}. Mức dung lượng này phải được chia sẻ đồng thời cho hệ điều hành, các ngăn xếp tác vụ (task stacks), bộ đệm truyền thông không dây và vùng nhớ cho mạng nơ-ron. Do đó, các mô hình học sâu buộc phải được nén tối ưu để tránh lỗi tràn bộ nhớ (Out of Memory) \cite{ref_1_11}.

Thứ hai là ràng buộc về độ trễ thực thi (Inference Latency). Trong bài toán giám sát máy giặt, tín hiệu rung động được lấy mẫu với tần số cao lên tới 1000 Hz. Điều này đòi hỏi mô hình phải hoàn thành quá trình suy luận (forward pass) trong một cửa sổ thời gian hẹp để đảm bảo tính thời gian thực, tránh gây thắt cổ chai dữ liệu. Cuối cùng là ràng buộc về năng lượng và năng lực xử lý; mô hình AI không được chiếm dụng quá nhiều chu kỳ CPU để đảm bảo hệ thống có thể duy trì đồng thời các tác vụ quản lý cảm biến và giao tiếp Wi-Fi/MQTT \cite{ref_1_1}.

\subsection{Kỹ thuật lượng tử hóa mô hình (Quantization INT8)}

Để vượt qua các giới hạn khắt khe về bộ nhớ và năng lực tính toán, lượng tử hóa nguyên 8-bit (INT8 Quantization) là kỹ thuật nén mô hình cốt lõi được áp dụng. Kỹ thuật này chuyển đổi các trọng số (weights) và giá trị kích hoạt (activations) từ định dạng số thực dấu phẩy động 32-bit (Float32) sang số nguyên 8-bit. Quá trình ánh xạ tuyến tính từ miền số thực $r$ sang miền số nguyên $q$ được mô tả qua Phương trình \ref{eq:quantization}:

\begin{equation}
\label{eq:quantization}
r = S(q - Z)
\end{equation}

Trong đó, $S$ (Scale) là số thực dương đại diện cho tỷ lệ thu phóng, và $Z$ (Zero-point) là giá trị số nguyên dùng để dịch tâm, tương ứng với giá trị 0 trong miền số thực \cite{ref_1_14}. Lượng tử hóa INT8 giúp giảm kích thước lưu trữ của mô hình đi khoảng 4 lần so với Float32, đồng thời giảm lượng RAM tĩnh tiêu thụ trong quá trình suy luận (inference) theo tỷ lệ tương ứng \cite{ref_1_1, ref_1_14}. Các nghiên cứu đã chỉ ra rằng khi kết hợp với phương pháp lượng tử hóa sau huấn luyện (Post-Training Quantization), mức suy giảm độ chính xác của mạng thường được kiểm soát dưới 1\% \cite{ref_1_14}. Hơn nữa, với kiến trúc lõi LX7 của ESP32-S3, dữ liệu kiểu INT8 cho phép vi điều khiển khai thác tối đa tập lệnh vector (Vector Instructions) chuyên dụng. Điều này giúp thực thi nhiều phép toán nhân-cộng tích lũy (MACs) song song, từ đó hệ thống có thể đạt được độ trễ suy luận khoảng 6,4 ms trên kiến trúc XIAO ESP32-S3 \cite{ref_1_4, ref_1_15}.

\subsection{Framework TensorFlow Lite for Microcontrollers}

Hệ thống sử dụng TensorFlow Lite for Microcontrollers (TFLite Micro) để nhúng trực tiếp mô hình tự mã hóa lượng tử hóa xuống vi điều khiển. TFLite Micro được thiết kế tối ưu cho các thiết bị biên, loại bỏ sự phụ thuộc vào hệ điều hành tiêu chuẩn và không yêu cầu cơ chế cấp phát bộ nhớ động. Việc không sử dụng các lệnh \texttt{malloc()} hay \texttt{new} trong quá trình chạy mô hình giúp ngăn ngừa rủi ro phân mảnh bộ nhớ, đảm bảo độ tin cậy cho hệ thống khi vận hành liên tục dài ngày \cite{ref_1_13}.

Framework này hoạt động bằng cách biên dịch mô hình thành một mảng byte không đổi (const byte array) lưu trên bộ nhớ Flash. Sau đó, một bộ thông dịch (Interpreter) sẽ điều phối quá trình tính toán thông qua một vùng đệm cấp phát sẵn duy nhất gọi là vùng đệm tensor (Tensor Arena) \cite{ref_1_13}. Việc tích hợp TFLite Micro trên hệ điều hành thời gian thực FreeRTOS cho phép hệ thống thực thi đa luồng độc lập. Dù tác vụ suy luận AI được gán mức ưu tiên thấp hơn các tác vụ thu thập dữ liệu phần cứng, kiến trúc đa nhân vẫn đảm bảo quá trình này đáp ứng tốt các yêu cầu về thời gian thực. Để mô hình TFLite Micro hoạt động chính xác, dữ liệu đầu vào cần đi qua một luồng xử lý tín hiệu số tối ưu, nội dung này sẽ được trình bày trong mục tiếp theo.

% --- Tài liệu tham khảo trích dẫn trong mục này ---
% \bibitem{ref_1_1} P. P. Ray, "A review on TinyML: State-of-the-art and prospects," 2022.
% \bibitem{ref_1_4} X. Wang et al., "Empowering Edge Intelligence: A Comprehensive Survey," 2025.
% \bibitem{ref_1_11} T. S. Ajani et al., "An Overview of Machine Learning within Embedded and Mobile Devices--Optimizations and Applications," Sensors, 2021.
% \bibitem{ref_1_13} Pete Warden and Daniel Situnayake, "TinyML: Machine Learning with TensorFlow Lite...", 2019.
% \bibitem{ref_1_14} B. Jacob et al., "Quantization and Training of Neural Networks...", CVPR 2018.
% \bibitem{ref_1_15} Espressif Systems, "ESP32-S3 Series Datasheet," Version 2.2, 2026.
